\section{Conclusion}
\label{sec:conclusion}

A validity rate for prompt-enforced structured output does not tell a
practitioner what failed or what a fix will change. Diagnosing a failure
requires keeping both the visible formatting failure and the schema
violation it may conceal; in this corpus, dropping either layer changes
the conclusion an engineer would draw about a repair, a schema change, or
a prompt change. Across
nine model configurations, those with similar rates fail
in different ways (RQ1); the dominant category moves when the binding constraints
change and differs between task domains for the same model (RQ2); and the
most common failure, markdown wrapping, conceals a length or pattern
violation often enough that reading primary labels alone misreads a
prompt ablation and misattributes what a repair loop leaves unrepaired
(RQ3). In both schema interventions we ran, whether a relaxation removed
a failure tracked the sub-cause beneath the label rather than the label
itself. Of the four practices we derive,
three are supported by these data; the fourth, building retry feedback
from the fence-stripped validation, follows from them but is untested. We release the classifier, raw
completions and annotation record so the corpus can be relabelled under
other definitions; the distribution over categories will move as models
change, and the vocabulary is meant to outlast it.
